\section{Tối ưu quy hoạch động}
\begingroup\small
\begin{multicols}{2}
\subsection{Convex Hull Trick}
Nếu chuyển trạng thái có dạng
\[
dp_i=\min_j\{m_jx_i+b_j\},
\]
mỗi $j$ là một đường thẳng $y=m_jx+b_j$; bài toán là
thêm đường và hỏi giá trị nhỏ nhất tại $x_i$.
Nếu độ dốc được thêm theo thứ tự \emph{giảm nghiêm ngặt}
và các truy vấn $x$ \emph{không giảm}, giữ bao dưới bằng deque/vector.
Với ba đường liên tiếp $a,b,c$ có $m_a>m_b>m_c$, đường $b$
bị bỏ nếu giao điểm $a,b$ không nằm trước giao điểm $b,c$:
\[
(b_b-b_a)(m_b-m_c)\ge
(b_c-b_b)(m_a-m_b).
\]
Đây là so sánh chéo chính xác bằng số nguyên, không dùng số thực.
Các tích của hiệu hệ số phải vừa \texttt{\_\_int128};
kiểu này chưa đủ nếu các hiệu có thể trải hết miền 64 bit.
Khi hai đường có cùng độ dốc, chỉ giữ đường có tung độ gốc nhỏ hơn.
Truy vấn tăng $x$ dùng con trỏ tiến; mỗi đường được thêm/xóa
không quá một lần, tổng $O(n)$ sau khi đã có thứ tự nói trên.
Phải thêm ít nhất một đường trước khi hỏi.
Nếu độ dốc hoặc truy vấn không theo thứ tự, không dùng mẫu này
mà đổi sang Li Chao hoặc cấu trúc khác.

\paragraph{Mẫu CHT đơn điệu (C++17).}
\cppsnippet{code/dp_optimization/cht.hpp}

\subsection{Li Chao tree}
Với miền truy vấn nguyên cố định $[L,R]$, Li Chao tree giữ ở
mỗi nút một đường tối ưu tại điểm giữa. Khi thêm đường mới,
so sánh ở đầu trái và điểm giữa, giữ đường tốt hơn tại giữa,
rồi đẩy đường còn lại xuống đúng một nửa có thể thắng.
Mỗi thao tác thêm/hỏi tốn $O(\log(R-L+1))$ và không cần
thứ tự độ dốc hay truy vấn. Mẫu dưới hỗ trợ \emph{min}, thêm
cả đường trên toàn miền; phải hỏi $x\in[L,R]$.
Giá trị hàm và đáp án dùng \texttt{\_\_int128}; vẫn cần kiểm
tra giới hạn tích $mx+b$.

\paragraph{Mẫu Li Chao động (C++17).}
\cppsnippet{code/dp_optimization/li_chao.hpp}

\subsection{Các tối ưu khác}
Với
$D_k(j)=\min_{i<j}\{D_{k-1}(i)+C(i,j)\}$,
nếu chỉ số tối ưu $\operatorname{opt}_k(j)$ không giảm theo $j$,
có thể tính một hàng bằng chia để trị: tính điểm giữa với
$i$ trong khoảng tối ưu cho phép, rồi thu hẹp khoảng cho
hai nửa. Chi phí $O(n\log n)$ lần gọi $C$ mỗi hàng
(nếu mỗi lần gọi $C$ là $O(1)$); không tự suy ra tính đơn
điệu chỉ từ vài ví dụ. Bất đẳng thức tứ giác/ma trận Monge
là một điều kiện đủ thường gặp. Hàm mẫu giả thiết
$D_{k-1}(i)$ đã biết và tính $j=1,\ldots,n$;
\texttt{prev[0]} là trạng thái trước vị trí đầu tiên. Giá trị vô cực
được mã hóa bằng \texttt{LLONG\_MAX/4}, mọi tổng hữu hạn phải
nằm trong $(-\mathrm{INF},\mathrm{INF})$.
Khi hòa chọn chỉ số nhỏ nhất, kể cả khi mọi ứng viên đều vô cực;
tính đơn điệu phải đúng với quy ước này.
\cppsnippet{code/dp_optimization/dnc.hpp}

Knuth áp dụng cho một dạng khác:
$D(l,r)=\min_{l\le k<r}\{D(l,k)+D(k+1,r)+W(l,r)\}$.
Nếu đã chứng minh được
$\operatorname{opt}(l,r-1)\le\operatorname{opt}(l,r)
\le\operatorname{opt}(l+1,r)$, chỉ thử $k$ trong khoảng
đó và giảm từ $O(n^3)$ xuống $O(n^2)$.
Điều kiện này không đồng nghĩa với điều kiện của chia để trị.
Khi chưa chứng minh được miền tối ưu, dùng DP gốc hoặc
kiểm chứng giả thiết trên dữ liệu nhỏ trước khi tối ưu.
\end{multicols}
\endgroup
